\chapter{The importance of an architecture}
\label{ch:prb}
\doublespacing
\vspace{-2.5cm}

\section{Foreword}
\label{sec:prb_foreword}
The development of an autonomous system is a process that involves various stages. Each stage is typically characterized by defining and solving issues related to different aspects of the system, ranging from the tasks it will be required to perform, to its construction and programming, and even to how the data it collects and that affect its operation should be presented to potential operators. Generalizing from a specific application case, most issues can be classified into a few high-level categories. Although it may seem counterintuitive from the division that follows, these categories are not always addressed sequentially, nor are they entirely independent of each other. To address situations like these, various project management techniques have been developed in recent times, aimed at minimizing the risk that an error made in an earlier phase could compromise subsequent stages, potentially undermining the work completed and extending development timeframes. In light of the observations made in \Cref{ch:intro}, it is clear that, with the increasing complexity of modern automated systems and of the tasks they must perform, these topics are more relevant than ever. In particular, new tools made available to developers and designers will need to reflect these specificities, simplify the management of typical situations where possible, and highlight potential critical issues to aid in problem-solving. Practical experience shows that, despite the variety of cases and applications, the challenges faced are almost always, if not identical, very similar, as are the strategies to address them.

The aim of this chapter is to briefly present the issues typically encountered when designing an autonomous system, highlighting their unique characteristics without focusing on a particular application area, and then to examine some recent hardware and software solutions that enable the optimization of design, development, and testing activities. The discussion starts at a relatively high level of abstraction but becomes progressively more specific, forming the guiding thread that connects all subsequent chapters, finally highlighting the necessity of a well-structured software architecture to address the issues raised.
